% ============ 第5章 §5.1 Schwartz 函数 (原书 pp.98-103) ============

这里, 作为深入考察至今为止出现的数学变换, 特别是 Hilbert 变换等的工具, 引入 Schwartz 函数空间与 Schwartz 分布空间.

\section{Schwartz 函数}\label{sec:5.1}
\index{0symboldx@$\mathrm{d}x'$}

\index{G@Gaussian}

\index{S@Schwartz 函数空间 $S(\mathbb{R}^{n})$}








在处理函数的性质时, 以 $L^{p}(\mathbb{R}^{n})$ 为代表的积分值等量的方面, 与以有界连续函数全体构成的 Banach 空间 $\mathrm{BC}(\mathbb{R}^{n})$ 为代表的连续性等质的方面, 都是重要的. 巧妙地组合这些性质, 就能进一步分析函数的性质. 兼备两者的函数是最理想的, 但作为这样的典型函数例子, 有具有紧支集的无限次可微函数全体构成的线性空间 $C_{\mathrm{c}}^{\infty}(\mathbb{R}^{n})$. 然而, $C_{\mathrm{c}}^{\infty}(\mathbb{R}^{n})$ 与数学中出现的变换, 例如积分变换等的相性不好. 这是因为 $C_{\mathrm{c}}^{\infty}(\mathbb{R}^{n})$ 的函数经过变换后, 往往不再是 $C_{\mathrm{c}}^{\infty}(\mathbb{R}^{n})$ 的函数. 因此, 考察比 $C_{\mathrm{c}}^{\infty}(\mathbb{R}^{n})$ 更大, 但被认为保持良好性质的函数空间. 令 $\mathbb{N}_{0} = \{0, 1, \ldots\}$, 用 $\mathbb{N}_{0}^{n}$ 表示多重指标全体构成的集合. 即,
\[
\mathbb{N}_{0}^{n} = \left\{(\alpha_{1}, \alpha_{2}, \ldots, \alpha_{n}) : \alpha_{j} \in \mathbb{N}_{0}, j = 1, 2, \ldots, n\right\}.
\]
对多重指标 $\beta = (\beta_{1}, \beta_{2}, \ldots, \beta_{n}) \in \mathbb{N}_{0}^{n}$ 与 $f \in C^{\infty}(\mathbb{R}^{n})$, 定义
\[
|\beta| = \sum_{j=1}^{n}\beta_{j} \quad (\beta \text{ 的长度}),
\]
\[
x^{\beta} \equiv x_{1}^{\beta_{1}}x_{2}^{\beta_{2}}\cdots x_{n}^{\beta_{n}} \quad (x \text{ 的 } \beta \text{ 次乘方}),
\]
\[
\partial^{\beta}f(x) \equiv \left(\frac{\partial}{\partial x_{1}}\right)^{\beta_{1}}\left(\frac{\partial}{\partial x_{2}}\right)^{\beta_{2}}\cdots\left(\frac{\partial}{\partial x_{n}}\right)^{\beta_{n}}f(x) \quad (f \text{ 的 } \beta \text{ 次偏导数}).
\]
又, 设 $\alpha = (\alpha_{1}, \alpha_{2}, \ldots, \alpha_{n})$, $\beta = (\beta_{1}, \beta_{2}, \ldots, \beta_{n})$ 为多重指标时, $\alpha \geq \beta$ 是指各分量满足相同的大小关系 $\alpha_{j} \geq \beta_{j}$. 并且, 此时用通常的二项式系数 ${}_{n}\mathrm{C}_{r}$ 定义一般的二项式系数 ${}_{\alpha}\mathrm{C}_{\beta}$:
\[
{}_{\alpha}\mathrm{C}_{\beta} = \prod_{j=1}^{n}{}_{\alpha_{j}}\mathrm{C}_{\beta_{j}}.
\]

\begin{definition}[Schwartz 函数空间 $\mathcal{S}(\mathbb{R}^{n})$]\label{def:5.1}
$\mathcal{S}(\mathbb{R}^{n})$ 是如下定义的函数集合. 对 $\alpha, \beta \in \mathbb{N}_{0}^{n}$, 暂且令
\begin{equation}\tag{5.1}
\varphi_{(\alpha,\beta)}(x) \equiv x^{\alpha}\partial^{\beta}\varphi(x) \quad (x \in \mathbb{R}^{n}),
\label{eq:5.1}
\end{equation}
将 $\mathcal{S}(\mathbb{R}^{n})$ 定义为对所有 $\alpha, \beta \in \mathbb{N}_{0}^{n}$ 满足 $\varphi_{(\alpha,\beta)} \in L^{\infty}(\mathbb{R}^{n})$ 的 $\varphi \in C^{\infty}(\mathbb{R}^{n})$ 全体构成的集合, $\mathcal{S}(\mathbb{R}^{n})$ 的元称为 Schwartz 函数.
\end{definition}

对 $\varphi \in C^{\infty}(\mathbb{R}^{n})$, 令
\[
p_{\alpha,\beta}(\varphi) = \sup_{x \in \mathbb{R}^{n}}|\varphi_{(\alpha,\beta)}(x)|.
\]

\begin{example}\label{exa:5.1}
对 $x \in \mathbb{R}^{n}$, 由 $E(x) \equiv e^{-\frac{1}{2}|x|^{2}}$ 给出的函数 $E$ 称为 Gaussian, 属于 $\mathcal{S}(\mathbb{R}^{n})$. 又 $C_{\mathrm{c}}^{\infty}(\mathbb{R}^{n}) \subset \mathcal{S}(\mathbb{R}^{n})$. 因此, $C_{\mathrm{c}}^{\infty}(\mathbb{R}^{n})$ 是 $\mathcal{S}(\mathbb{R}^{n})$ 的真子集.
\end{example}

在下面的例子中, 使用多项式定理
\[
(A_{1} + A_{2} + \cdots + A_{n})^{k} = \sum_{k_{1}, k_{2}, \ldots, k_{n} \in \mathbb{N}_{0}, k_{1} + k_{2} + \cdots + k_{n} = k}\frac{k!}{k_{1}!k_{2}!\cdots k_{n}!}A_{1}^{k_{1}}A_{2}^{k_{2}}\cdots A_{n}^{k_{n}}.
\]

\begin{example}\label{exa:5.2}
设 $\varphi \in \mathcal{S}(\mathbb{R}^{n})$. 由上述注意,
\begin{align*}
(1+|x|)^{n+1} &\leq 2^{n+1}(1+|x|^{2})^{n+1} \\
&= 2^{n+1}\sum_{\substack{k_{0}, k_{1}, k_{2}, \ldots, k_{n} \in \mathbb{N}_{0} \\ k_{0}+k_{1}+k_{2}+\cdots+k_{n}=n+1}}\frac{(n+1)!|x^{2k}|}{k_{1}!k_{2}!\cdots k_{n}!k_{0}!}
\end{align*}
所以 $\Phi(x) \equiv (1+|x|)^{n+1}\varphi(x)$ 是有界的. 因此, $|\varphi(x)| \leq (1+|x|)^{-n-1}\|\Phi\|_{L^{\infty}}$, 由例~\ref{exa:2.1}, $\varphi \in L^{1}(\mathbb{R}^{n})$. 由于对任意多重指标 $\alpha, \beta$, $\varphi_{(\alpha,\beta)} \in \mathcal{S}(\mathbb{R}^{n})$, 所以 $\varphi_{(\alpha,\beta)} \in L^{1}(\mathbb{R}^{n})$.
\end{example}

\begin{example}\label{exa:5.3}
设 $f, g \in C^{\infty}(\mathbb{R}^{n})$. 由广义 Leibniz 法则,
\[
\partial^{\alpha}[f(x)g(x)] = \sum_{\beta \leq \alpha}{}_{\alpha}\mathrm{C}_{\beta}\partial^{\beta}f(x) \cdot \partial^{\alpha-\beta}g(x)
\]
成立. 因此, 设 $f, g \in \mathcal{S}(\mathbb{R}^{n})$, 则 $f \cdot g \in \mathcal{S}(\mathbb{R}^{n})$.
\end{example}

\begin{example}\label{exa:5.4}
设 $M > 0$. 对满足 $|h(x)| + |\nabla h(x)| \lesssim (1+|x|)^{M}$, $x \in \mathbb{R}^{n}$ 的 $C^{1}$ 级函数 $h$ 与 $f \in \mathcal{S}(\mathbb{R}^{n})$, 分部积分公式
\[
\int_{\mathbb{R}^{n}}\partial_{j}f(x) \cdot h(x)\dd x = -\int_{\mathbb{R}^{n}}f(x) \cdot \partial_{j}h(x)\dd x \quad (j = 1, 2, \ldots, n)
\]
成立. 由例~\ref{exa:2.1}, 注意两边都是有限的.

实际上, 考虑乘积的微分, 只需证明
\begin{equation}\tag{5.2}
\int_{\mathbb{R}^{n}}\partial_{j}[f(x) \cdot h(x)]\dd x = 0
\label{eq:5.2}
\end{equation}
即可.

由于具有对称性, 以 $j = n$ 的情形证明. 记 $x' = (x_{1}, x_{2}, \ldots, x_{n-1})$, $\dd x' = \dd x_{1}\dd x_{2}\cdots\dd x_{n-1}$. 由例~\ref{exa:2.1} 与 Fubini 定理,
\[
\int_{\mathbb{R}^{n}}\partial_{n}[f(x) \cdot h(x)]\dd x = \int_{\mathbb{R}^{n-1}}\left(\int_{\mathbb{R}}\partial_{n}[f(x) \cdot h(x)]\dd x_{n}\right)\dd x'
\]
可以这样分解. 固定 $x' \in \mathbb{R}^{n-1}$. 将微积分基本定理
\[
\int_{-R}^{R}\varphi'(t)\dd t = \varphi(R) - \varphi(-R)
\]
应用于 $\varphi(t) = f(x', t)h(x', t)$,
\begin{align*}
\int_{\mathbb{R}}\partial_{n}[f(x) \cdot h(x)]\dd x_{n} &= \lim_{R \to \infty}\int_{-R}^{R}\partial_{n}[f(x) \cdot h(x)]\dd x_{n} \\
&= \lim_{R \to \infty}(f(x', R)h(x', R) - f(x', -R)h(x', -R)) \\
&= 0
\end{align*}
成立. 由以上, 可知式~\eqref{eq:5.2} 成立.
\end{example}

作为集合的 $\mathcal{S}(\mathbb{R}^{n})$ 虽然如此定义, 但可以给 $\mathcal{S}(\mathbb{R}^{n})$ 赋予距离. 即, 当 $f, g \in \mathcal{S}(\mathbb{R}^{n})$ 时, 通过定义 f, g 的距离 $d(f, g)$ 为
\[
d(f, g) = \sum_{\alpha, \beta \in \mathbb{N}_{0}^{n}}\frac{p_{\alpha,\beta}(f-g)}{2^{|\alpha|+|\beta|}p_{\alpha,\beta}(f-g)+1},
\]
$(\mathcal{S}(\mathbb{R}^{n}), d)$ 是距离空间. 由
\[
\frac{p_{\alpha,\beta}(f-g)}{2^{|\alpha|+|\beta|}p_{\alpha,\beta}(f-g)+1} \leq \frac{1}{2^{|\alpha|+|\beta|}}, \quad \sum_{\alpha, \beta \in \mathbb{N}_{0}^{n}}\frac{1}{2^{|\alpha|+|\beta|}} < \infty
\]
可以使用 Lebesgue 收敛定理, 从而可以说出下面的事.

\begin{theorem}\label{thm:5.1}
$\{f_{j}\}_{j=1}^{\infty} \subset \mathcal{S}(\mathbb{R}^{n})$ 收敛于 $f \in \mathcal{S}(\mathbb{R}^{n})$ 的必要充分条件是, 对所有 $\alpha, \beta \in \mathbb{N}_{0}^{n}$,
\[
\lim_{j \to \infty}p_{\alpha,\beta}(f_{j}-f) = 0
\]
成立.
\end{theorem}

对 $\varphi \in \mathcal{S}(\mathbb{R}^{n})$, 令
\[
p_{N}(\varphi) \equiv \sum_{|\beta| \leq N}\sup_{x \in \mathbb{R}^{n}}(1+|x|)^{N}|\partial^{\beta}\varphi(x)|.
\]
也可以令
\[
p_{N}(\varphi) \equiv \sum_{|\alpha|, |\beta| \leq N}\sup_{x \in \mathbb{R}^{n}}|x^{\alpha}\partial^{\beta}\varphi(x)|,
\]
但暂且为统一起见采用前者. 虽然无论采用哪一个都是相同的, 但作为作者个人的感觉, 有着想在可能的定义中采用给出最大值的定义作为定义的心理. 即, 基于
\[
\sum_{|\alpha|, |\beta| \leq N}\sup_{x \in \mathbb{R}^{n}}|x^{\alpha}\partial^{\beta}\varphi(x)| \leq (N+1)^{n}\sum_{|\beta| \leq N}\sup_{x \in \mathbb{R}^{n}}(1+|x|)^{N}|\partial^{\beta}\varphi(x)|
\]
想采用前者作为定义.

\begin{example}\label{exa:5.5}
$p_{N}, p_{\alpha,\beta}$ 等被称为泛函, 但引入 $p_{N}, p_{\alpha,\beta}$ 的目的, 并不是具体计算各 $\varphi$ 的值 $p_{N}(\varphi), p_{\alpha,\beta}(\varphi)$, 而是为了讨论这些值的有限性, 或者用 $p_{N}, p_{\alpha,\beta}$ 讨论某种函数列的收敛性. 因此, 将这些自由替换来考虑也可以. 立足于这一观点, 可知除这些之外还可能有各种各样的泛函. 举例来看. 对 $\varphi \in \mathcal{S}(\mathbb{R}^{n})$, 定义
\[
\begin{aligned}
p_{N}(\varphi) = p_{N}^{1}(\varphi) &= \sup_{\beta \in \mathbb{N}_{0}^{n}, |\beta| \leq N}\sup_{x \in \mathbb{R}^{n}}(1+|x|)^{N}|\partial^{\beta}\varphi(x)|, \\
p_{N}^{2}(\varphi) &= \sup_{\beta \in \mathbb{N}_{0}^{n}, |\beta| \leq N}\sup_{x \in \mathbb{R}^{n}}(1+|x|^{2})^{N}|\partial^{\beta}\varphi(x)|
\end{aligned}
\]
由于 $1+|x| \leq 2(1+|x|^{2})$, $1+|x|^{2} \leq (1+|x|)^{2}$,
\[
p_{N}^{1}(\varphi) \leq 2^{N}p_{N}^{2}(\varphi) \leq 2^{N}p_{2N}^{1}(\varphi)
\]
成立. 因此, $\{p_{N}^{1}\}_{N=1}^{\infty}$, $\{p_{N}^{2}\}_{N=1}^{\infty}$ 无论将哪一个作为 $\{p_{N}\}_{N=1}^{\infty}$ 的定义都可以.
\end{example}

\begin{example}\label{exa:5.6}
设 $O \subset \mathcal{S}(\mathbb{R}^{n})$ 是包含 0 函数的开集. 此时存在某个 $L > 0$, 使得
\[
\{\varphi \in \mathcal{S}(\mathbb{R}^{n}) : Lp_{L}(\varphi) < 1\} \subset O
\]
成立. 实际上, 由于开集由 d 生成, 存在某个 $r > 0$, 使得
\[
\{\varphi \in \mathcal{S}(\mathbb{R}^{n}) : d(\varphi, 0) < r\} \subset O
\]
成立. 取自然数 N 使得
\[
\sum_{\alpha, \beta \in \mathbb{N}_{0}^{n}}\frac{\chi_{(N,\infty)}(\max(|\alpha|, |\beta|))}{2^{|\alpha|+|\beta|}} < \frac{r}{2}.
\]
对这样的 N, 选取 $L \geq N$ 使得
\[
L^{-1}\sharp\{\alpha, \beta \in \mathbb{N}_{0}^{n} : \max(|\alpha|, |\beta|) \leq N\} < \frac{r}{2}.
\]
于是, 若 $Lp_{L}(\varphi) < 1$, 则
\begin{align*}
d(\varphi, 0) &= \sum_{\alpha, \beta \in \mathbb{N}_{0}^{n}}\frac{p_{\alpha,\beta}(\varphi)\chi_{(N,\infty)}(\max(|\alpha|, |\beta|))}{2^{|\alpha|+|\beta|}p_{\alpha,\beta}(\varphi)+1} \\
&\quad + \sum_{\alpha, \beta \in \mathbb{N}_{0}^{n}}\frac{p_{\alpha,\beta}(\varphi)\chi_{[0,N]}(\max(|\alpha|, |\beta|))}{2^{|\alpha|+|\beta|}p_{\alpha,\beta}(\varphi)+1} \\
&\leq \sum_{\alpha, \beta \in \mathbb{N}_{0}^{n}}\frac{\chi_{(N,\infty)}(\max(|\alpha|, |\beta|))}{2^{|\alpha|+|\beta|}} \\
&\quad + L^{-1}\sharp\{\alpha, \beta \in \mathbb{N}_{0}^{n} : \max(|\alpha|, |\beta|) \leq N\} \\
&< \frac{r}{2} + \frac{r}{2} = r
\end{align*}
所以
\[
\{\varphi \in \mathcal{S}(\mathbb{R}^{n}) : Lp_{L}(\varphi) < 1\} \subset \{\varphi \in \mathcal{S}(\mathbb{R}^{n}) : d(\varphi, 0) < r\} \subset O
\]
成立.
\end{example}

由下面的例子可知, $C_{\mathrm{c}}^{\infty}(\mathbb{R}^{n})$ 在 $\mathcal{S}(\mathbb{R}^{n})$ 中是稠密的.

\begin{example}\label{exa:5.7}
设 $f \in \mathcal{S}(\mathbb{R}^{n})$, $\psi \in \mathcal{S}(\mathbb{R}^{n})$ 满足 $\chi_{B(1)} \leq \psi \leq \chi_{B(2)}$. 此时, 在 $\mathcal{S}(\mathbb{R}^{n})$ 的拓扑中, 当 $j \to \infty$ 时, $\psi(2^{-j}\cdot)f \to f$ 成立. 实际上, 取 $\alpha, \beta \in \mathbb{N}_{0}^{n}$, 基于定理~\ref{thm:5.1} 证明
\[
\lim_{j \to \infty}p_{\alpha,\beta}(\psi(2^{-j}\cdot)f - f) = 0.
\]
由 Leibniz 法则,
\begin{align*}
&x^{\alpha}(\partial^{\beta}(\psi(2^{-j}x)f(x) - f(x))) \\
&= \sum_{0 \neq \gamma \leq \beta}2^{-j|\gamma|}{}_{\beta}\mathrm{C}_{\gamma}x^{\alpha}\partial^{\gamma}\psi(2^{-j}x) \cdot \partial^{\beta-\gamma}f(x) \\
&\quad + x^{\alpha}(\psi(2^{-j}x)\partial^{\beta}f(x) - \partial^{\beta}f(x))
\end{align*}
这里, 使用
\begin{align*}
&\sup_{x \in \mathbb{R}^{n}}|x^{\alpha}\partial^{\gamma}\psi(2^{-j}x) \cdot \partial^{\beta-\gamma}f(x)| \leq \|\partial^{\gamma}\psi\|_{L^{\infty}}\sup_{x \in \mathbb{R}^{n}}|x^{\alpha}\partial^{\beta-\gamma}f(x)| < \infty, \\
&\sup_{x \in \mathbb{R}^{n}}|x^{\alpha}[\psi(2^{-j}x)\partial^{\beta}f(x) - \partial^{\beta}f(x)]| \\
&= \sup_{x \in \mathbb{R}^{n}\setminus B(2^{j})}|x^{\alpha}[\psi(2^{-j}x)\partial^{\beta}f(x) - \partial^{\beta}f(x)]| \\
&\leq \sup_{x \in \mathbb{R}^{n}\setminus B(2^{j})}|x^{\alpha}\partial^{\beta}f(x)| \\
&\leq 2^{-j}\sup_{x \in \mathbb{R}^{n}\setminus B(2^{j})}|x| \cdot |x^{\alpha}\partial^{\beta}f(x)|
\end{align*}
得到
\[
p_{\alpha,\beta}(\psi(2^{-j}\cdot)f - f) = \sup_{x \in \mathbb{R}^{n}}|x^{\alpha}(\partial^{\beta}(\psi(2^{-j}x)f(x) - f(x)))| = O(2^{-j}).
\]
因此 $\lim_{j \to \infty}p_{\alpha,\beta}(\psi(2^{-j}\cdot)f - f) = 0$.
\end{example}

\begin{exercise}\label{exer:5.1}
设 $A, B, C > 0$. 证明
\[
\frac{B}{AB+1} + \frac{C}{AC+1} \geq \frac{B+C}{A(B+C)+1}.
\]
【提示】当 $t \geq 0$ 时, 令 $\varphi(t) = \frac{t}{t+1}$, 则问题中的不等式可以改写为
\[
\varphi(AB + AC) \leq \varphi(AB) + \varphi(AC)
\]
($\iff \varphi(AB + AC) + \varphi(0) \leq \varphi(AB) + \varphi(AC)$).
利用
\[
\varphi(AB) - \varphi(AB + AC) = -\int_{AB}^{AB+AC}\varphi'(t)\dd t,
\]
\[
\varphi(AC) - \varphi(0) = \int_{0}^{AC}\varphi'(t)\dd t,
\]
证明起来很容易.
\end{exercise}

\begin{exercise}\label{exer:5.2}
证明 $(\mathcal{S}(\mathbb{R}^{n}), d)$ 是完备的距离空间.
\end{exercise}

\begin{exercise}\label{exer:5.3}
设 $1 \leq p \leq \infty$. 利用 $\mathcal{S}(\mathbb{R}^{n})$ 的任意元 $\varphi$ 满足
\[
q(\varphi) = \sup_{x \in \mathbb{R}^{n}}|\varphi(x)|, \quad r_{j}(\varphi) = \sup_{x \in \mathbb{R}^{n}}|x_{j}|^{n+1}|\varphi(x)| < \infty
\]
这一点, 回答下面的问题.

(1) 必要时使用
\[
(1+|x|) \geq \prod_{j=1}^{n}\sqrt[n]{1+|x_{j}|},
\]
证明 $f(x) = \frac{1}{(1+|x|)^{n+1}} \in L^{p}(\mathbb{R}^{n})$.

(2) 设 $g \in \mathcal{S}(\mathbb{R}^{n})$. 必要时使用
\[
\begin{aligned}
(1+|x|)^{n+1} &\leq \left(1+\sum_{j=1}^{n}|x_{j}|\right)^{n+1} \\
&\leq (n+1)^{n+1} + (n+1)^{n+1}\sum_{j=1}^{n}|x_{j}|^{n+1},
\end{aligned}
\]
证明 $M = \sup_{x \in \mathbb{R}^{n}}(1+|x|)^{n+1}|g(x)| < \infty$.

(3) 证明 $\mathcal{S}(\mathbb{R}^{n}) \subset L^{p}(\mathbb{R}^{n})$.
\end{exercise}

\begin{exercise}\label{exer:5.4}
当 $f, g \in \mathcal{S}(\mathbb{R}^{n})$ 时, 证明 $f \cdot g \in \mathcal{S}(\mathbb{R}^{n})$. (补充上述例5.3 的细节.)
\end{exercise}
